% Paper-ready Results section for ReEntry-TAP
% Generated: 2026-07-04

\section{Experiments and Results}
\label{sec:experiments}

\subsection{Evaluation protocol}
\label{sec:evaluation_protocol}

We evaluate ReEntry on standard TAP-Vid datasets and re-entry-focused diagnostic settings. For standard TAP-Vid metrics, we report Average Jaccard (AJ), Occlusion Accuracy (OA), and average position accuracy ($\delta_{\mathrm{avg}}$), computed using the official TAP-Vid metric formulation or an equivalent local port. Since ReEntry targets a sparse failure mode that occurs after a point reappears following occlusion or disappearance, we additionally report AJ\_RD as a re-entry diagnostic metric. AJ\_RD is not an official leaderboard metric; it isolates the post-reappearance subset that can be diluted by full-trajectory averages.

Unless otherwise stated, all results are local evaluations under the specified protocol, not official server or leaderboard submissions. Table~\ref{tab:evaluation_protocols} summarizes the evaluation settings used in this work.

\subsection{DAVIS first-query evaluation}
\label{sec:davis_first_input}

Table~\ref{tab:davis_first_input} reports results on TAPVid-DAVIS under a first-query/input-resolution official-style local evaluation. This setting is the most favorable official-style DAVIS protocol for the current ReEntry family. Compared with the CoTracker3 offline base, ReEntry improves both standard TAP-Vid metrics and the re-entry diagnostic.

The strongest ReEntry variant in this setting is V24-DINOScore. It improves AJ from $62.6566$ to $64.8439$, OA from $88.1487$ to $91.8851$, and AJ\_RD from $0.3142$ to $0.3549$. This corresponds to gains of $+2.1874$ AJ, $+3.7363$ OA, and $+0.0407$ AJ\_RD over the CoTracker3 offline base. The position metric $\delta_{\mathrm{avg}}$ remains unchanged because ReEntry modifies visibility decisions rather than coordinates.

Paired-video statistics support this improvement. V24-DINOScore improves AJ over the offline base by $+2.1874$ percentage points with a 95\% bootstrap confidence interval of $[+1.2582,+3.2487]$, and improves OA by $+3.7363$ percentage points with a confidence interval of $[+2.1673,+5.7423]$. V24 also slightly improves over V22Q, with $+0.1179$ AJ and $+0.1445$ OA.

TrackOn2 remains the strongest external baseline on DAVIS first/input, reaching $67.0406$ AJ and $92.0916$ OA. We therefore do not claim that ReEntry surpasses TrackOn2 on DAVIS. Instead, this result shows that ReEntry substantially improves the CoTracker3 offline base in the query-first setting while preserving the intended method scope.

\subsection{DAVIS strided/original evaluation}
\label{sec:davis_strided_original}

Table~\ref{tab:davis_strided_original} reports results on TAPVid-DAVIS under a stricter strided/original official-style local evaluation. This setting exposes an important protocol sensitivity. The default ReEntry variants improve AJ\_RD but reduce global AJ/OA. V1 increases AJ\_RD from $0.3870$ to $0.4144$, and V22Q increases AJ\_RD to $0.4136$. However, their standard AJ and OA decrease relative to the CoTracker3 offline base.

This behavior reflects an objective mismatch. ReEntry is designed to correct visibility lag after reappearance. Under strided evaluation, visibility predictions are scored over a broader set of query frames and time steps. Aggressive visibility recovery can therefore introduce false-positive visibility decisions that improve re-entry recovery while reducing global AJ/OA.

To test whether ReEntry can be made official-style safe, we run a threshold sweep and identify V25-safe with threshold $0.80$. This conservative operating point reaches $51.5648$ AJ and $92.2106$ OA, slightly above the offline base of $51.5385$ AJ and $92.1543$ OA, while retaining a small AJ\_RD gain from $0.3870$ to $0.3900$. Paired-video intervals cross zero, so this should not be interpreted as a strong leaderboard-style improvement. We instead treat it as an official-style safe operating point.

This result motivates future metric-aware selective ReEntry: rather than firing whenever re-entry recovery appears likely, the selector should also estimate the risk of harming standard AJ/OA.

\subsection{RGB-Stacking full50 local standard evaluation}
\label{sec:rgb_full50}

Table~\ref{tab:rgb_full50} reports results on the complete locally available TAPVid RGB-Stacking full50 evaluation, covering videos \texttt{rgb\_stacking\_000000} through \texttt{rgb\_stacking\_000049}. This is the strongest full local standard evidence for ReEntry on the RGB-Stacking benchmark.

The CoTracker3 offline base obtains $0.3617$ AJ\_RD, $79.9345$ AJ, and $91.6371$ OA. ReEntry V1 improves AJ\_RD to $0.4414$ and OA to $93.0758$, corresponding to $+0.0797$ AJ\_RD and $+1.4387$ OA. ReEntry V22Q and V24-DINOScore provide similar re-entry gains with slightly better AJ preservation than V1. V24-DINOScore reaches $0.4394$ AJ\_RD, $79.5974$ AJ, and $93.0389$ OA.

The trade-off is a small reduction in standard AJ. V24-DINOScore improves AJ\_RD by $+0.0777$ and OA by $+1.4018$, but reduces AJ by $-0.3371$. Paired-video statistics show that the AJ\_RD and OA improvements are stable: V24 improves AJ\_RD on 48 videos with no losses and 2 ties, and improves OA on 40 videos with 10 losses. This supports the central claim that ReEntry improves re-entry recovery and visibility stability on RGB-Stacking full50, while introducing a small AJ trade-off.

\subsection{Re-entry stress diagnostics}
\label{sec:stress_diagnostics}

Table~\ref{tab:fresh_stress} isolates the re-entry failure mode using RGB-Stacking fresh20--49 and synthetic translate/occluder stress variants. These settings are diagnostic rather than official benchmark submissions. They are designed to evaluate whether ReEntry remains useful when reappearance and occlusion-related failures are amplified.

Across natural, translate-L16, and occluder-L16 settings, V24-DINOScore consistently improves AJ\_RD and OA over the base. Averaged over the three settings, AJ\_RD improves from $0.4939$ to $0.5451$, a gain of $+0.0509$. OA improves from $91.0453$ to $92.4496$, a gain of $+1.4043$. AJ decreases from $77.4721$ to $77.1716$, a small trade-off of $-0.3006$.

These diagnostics confirm that the method addresses the intended failure mode: ReEntry improves post-reappearance recovery and visibility behavior, but aggressive recovery can slightly reduce standard AJ.

\subsection{Ablation and method positioning}
\label{sec:method_positioning}

The current ReEntry family should be interpreted as follows.

\textbf{V1 Learned ReEntry-VisCalibrator} is the AJ\_RD-oriented main method. It provides the strongest re-entry recovery on RGB-Stacking full50 and strong improvements over the CoTracker3 offline base in DAVIS first/input.

\textbf{V22Q Interval/Gate} is the stability-oriented extension. It slightly reduces aggressive visibility updates and improves standard AJ relative to V1 in several settings.

\textbf{V24-DINOScore} is an optional appearance micro-filter. It provides small standard-metric improvements over V22Q in DAVIS first/input and slightly better AJ preservation on RGB-Stacking full50, but it should not be presented as a universal replacement for V1/V22Q.

\textbf{V25-safe threshold $0.80$} is an official-style safe operating point for DAVIS strided/original. It preserves AJ/OA while retaining only a small AJ\_RD gain. It is useful for protocol-safe reporting, but not a strong leaderboard-style method.

\subsection{Summary of findings}
\label{sec:results_summary}

The experiments support four conclusions.

First, ReEntry substantially improves the CoTracker3 offline base under DAVIS first/input evaluation, improving AJ, OA, and AJ\_RD without changing coordinates.

Second, ReEntry improves re-entry recovery and OA on RGB-Stacking full50, the strongest full local standard evidence in the current evaluation suite, with a small standard-AJ trade-off.

Third, diagnostic stress evaluations confirm that the improvements are concentrated on the intended failure mode: post-reappearance recovery and visibility correction.

Fourth, official strided/original evaluation reveals a limitation. Default ReEntry variants improve AJ\_RD but can harm global AJ/OA. A conservative V25 threshold can make the method official-style safe, but the resulting gains are marginal. This motivates future official-metric-aware selective ReEntry.

Overall, ReEntry should be claimed as a re-entry failure-mode method, not as a universal TAP-Vid leaderboard booster. The strongest safe claim is that ReEntry improves re-entry recovery across standard TAP-Vid local evaluations and diagnostic stress settings; it improves DAVIS first-query AJ/OA/AJ\_RD and RGB-Stacking full50 AJ\_RD/OA, while stricter strided/original evaluation requires conservative thresholding to preserve official-style AJ/OA.
